\documentclass[11pt]{article}
\usepackage[letterpaper,margin=1.1in]{geometry}
\usepackage{mathptmx}
\usepackage[T1]{fontenc}
\usepackage{microtype}
\setlength{\parindent}{0pt}
\setlength{\parskip}{0.6em}
\newcommand{\sect}[1]{\vspace{0.9em}{\large\textbf{#1}}\par\vspace{-0.2em}}
\begin{document}

\noindent Linguistics 150C \hfill \textbf{Sample: preliminary sources and initial research question}\\[-0.6em]
\noindent\rule{\linewidth}{0.4pt}

\textit{This is an example of the format and level of detail expected for the November 6 submission. It is on a topic that is not part of the course, so please choose your own topic and sources rather than building on this one.}

\sect{Initial research question}
How has contact with Spanish, and the shift toward Spanish in daily life, affected the use and structure of Rapa Nui, the Polynesian language of Easter Island (Rapa Nui), since the island's annexation by Chile in 1888?

I may narrow this to one area as I read, either vocabulary borrowing or the situations in which each language is used.

\sect{Preliminary sources}
\textbf{1.} Du Feu, Veronica. 1996. \textit{Rapanui}. London: Routledge.

This is a descriptive grammar of Rapa Nui. I will use it to find out what the language looked like structurally (sounds, word order, vocabulary) so I can tell which features are inherited from Proto-Polynesian and which may have come from Spanish.

\textbf{2.} Makihara, Miki. 2005. Rapa Nui ways of speaking Spanish: Language shift and socialization on Easter Island. \textit{Language in Society} 34(5): 727--762.

This article looks at how Rapa Nui people use Spanish and Rapa Nui and how children are socialized into each. It directly addresses the part of my question about language shift and when each language is used.

\textbf{3.} Fischer, Steven Roger. 1997. \textit{Rongorongo: The Easter Island script}. Oxford: Clarendon Press.

This book covers the island's history and its script. I plan to use its historical background (contact, depopulation, and annexation) to place the linguistic changes in time, and to see how literacy relates to language use.

\sect{What I expect to learn}
I expect to find that Rapa Nui borrowed many words from Spanish, as many Pacific languages have from colonial languages, but that the larger change is in who speaks which language and where. I will look for evidence that Spanish is used in more settings than Rapa Nui, and I will try to find out whether Spanish has also changed Rapa Nui grammar, or only its vocabulary. I am not sure yet whether my sources will give enough detail on grammar, which may lead me to focus on vocabulary and language use.

\end{document}
